\documentclass[11pt,a4paper]{article}
\usepackage{fontspec}
\setmainfont{texgyrepagella}[Extension=.otf,UprightFont=*-regular,BoldFont=*-bold,ItalicFont=*-italic,BoldItalicFont=*-bolditalic]
\usepackage{unicode-math}
\setmathfont{texgyrepagella-math.otf}
\setsansfont{texgyreheros}[Extension=.otf,UprightFont=*-regular,BoldFont=*-bold,ItalicFont=*-italic,BoldItalicFont=*-bolditalic]
\setmonofont{lmmono10-regular.otf}[Scale=0.85]
\usepackage[margin=2.4cm,top=2.2cm,bottom=2.6cm]{geometry}
\usepackage{xcolor}
\definecolor{ink}{HTML}{1B2A41}
\definecolor{accent}{HTML}{B4442E}
\definecolor{soft}{HTML}{F5F1EA}
\definecolor{rule}{HTML}{D8CFC0}
\usepackage[most]{tcolorbox}
\usepackage{longtable,booktabs,array,calc,tabularx}
\usepackage{enumitem}
\setlist{itemsep=2pt,topsep=4pt}
\usepackage{titlesec}
\titleformat{\section}{\sffamily\Large\bfseries\color{ink}}{}{0pt}{}[\vspace{-4pt}{\color{rule}\rule{\linewidth}{0.6pt}}]
\titleformat{\subsection}{\sffamily\large\bfseries\color{ink}}{}{0pt}{}
\setcounter{secnumdepth}{0}
\titlespacing*{\section}{0pt}{18pt}{8pt}
\usepackage{fancyhdr}
\pagestyle{fancy}\fancyhf{}
\renewcommand{\headrulewidth}{0pt}
\fancyhead[L]{\sffamily\footnotesize\color{ink!60}Kourovka Notebook,
Problem 20.87}
\fancyhead[R]{\sffamily\footnotesize\color{ink!60}Preprint, October
2026}
\fancyfoot[C]{\sffamily\footnotesize\color{ink!60}\thepage}
\usepackage{hyperref}
\hypersetup{colorlinks=true,linkcolor=accent,urlcolor=accent,citecolor=accent,pdftitle={Finite
simple groups whose element orders have at most two prime
divisors},pdfauthor={AutoPasha (AI employee, autopasha.ru)}}
\usepackage{microtype}\emergencystretch=3em
\setlength{\parindent}{0pt}\setlength{\parskip}{5pt}
\AtBeginDocument{\providecommand{\square}{\mdwhtsquare}\providecommand{\blacksquare}{\mdblksquare}}
\providecommand{\tightlist}{\setlength{\itemsep}{1pt}\setlength{\parskip}{0pt}}
\newtcolorbox{answerbox}{enhanced,breakable,colback=soft,colframe=accent,boxrule=0.8pt,arc=2pt,left=10pt,right=10pt,top=8pt,bottom=8pt,
  borderline west={3pt}{0pt}{accent}}
\newtcolorbox{abstractbox}{enhanced,colback=white,colframe=rule,boxrule=0.5pt,arc=2pt,left=12pt,right=12pt,top=8pt,bottom=8pt}
\newtcolorbox{notebox}{enhanced,breakable,colback=ink!4,colframe=ink!4,arc=2pt,left=10pt,right=10pt}
\begin{document}
\thispagestyle{empty}
\begin{tcolorbox}[enhanced,colback=ink,colframe=ink,arc=0pt,left=16pt,right=16pt,top=14pt,bottom=14pt,width=\linewidth]
{\sffamily\color{white!70!ink}\small PREPRINT · KOUROVKA NOTEBOOK
20.87}\par\vspace{6pt}
{\sffamily\color{white}\LARGE\bfseries Finite simple groups whose
element orders have at most two prime divisors}\par\vspace{6pt}
{\color{white!85!ink}\large A solution of Problem 20.87 of the Kourovka
Notebook}\par\vspace{10pt}
{\sffamily\color{white!75!ink}\small AutoPasha (AI employee,
autopasha.ru) \quad·\quad October 2026}
\end{tcolorbox}
\vspace{6pt}
\begin{abstractbox}

\textbf{Abstract.} We answer Problem 20.87 of the Kourovka Notebook (W.
J. Shi): the non-abelian composition factors of finite groups in which
every element order is divisible by at most two primes are exactly the
finite non-abelian simple groups with this property, and we classify
them. Besides the alternating groups \(A_n\), \(5\le n\le 11\), and
thirteen sporadic groups, they are the groups \(L_2(q)\), \(L_3(q)\),
\(U_3(q)\), \(Sz(q)\), \({}^2G_2(q)\) satisfying explicit arithmetic
conditions on \(q\), and eighteen further groups of Lie type.

\emph{AI disclosure.} This solution was obtained by an AI system
(AutoPasha, an AI employee, autopasha.ru) and checked by independent
computer verification (GAP 4.14 and its Character Table Library). It has
not yet been reviewed by a human mathematician; the GAP scripts are
available from the authors.

\end{abstractbox}

\textbf{Problem 20.87.} \emph{What are the non-abelian composition
factors of finite groups in which the order of every element is
divisible by at most two primes?} (For one prime, i.e.~EPPO-groups, see
arXiv:2003.09445.)

Throughout, \(\omega(n)\) is the number of distinct prime divisors of a
positive integer \(n\) (\(\omega(1)=0\)), and we say that a finite group
\(X\) has \textbf{property (P2)} if \(\omega(|x|)\le 2\) for every
\(x\in X\), i.e.~no element order of \(X\) is divisible by three
distinct primes. \(p\) always denotes the defining characteristic of a
group of Lie type over \(\mathbb F_q\).

\subsection{0. The answer}\label{the-answer}

\begin{answerbox}

\textbf{Main Theorem.} A finite non-abelian simple group \(T\) is a
composition factor of a finite group with property (P2) if and only if
\(T\) itself has (P2), and this happens exactly when \(T\) is isomorphic
to one of the following groups.

\begin{enumerate}
\def\labelenumi{\arabic{enumi}.}
\tightlist
\item
  \textbf{Alternating groups:} \(A_n\), \(5\le n\le 11\).
\item
  \textbf{Sporadic groups:} \(M_{11}\), \(M_{12}\), \(M_{22}\),
  \(M_{23}\), \(M_{24}\), \(J_1\), \(J_2\), \(J_3\), \(HS\), \(He\),
  \(Ru\), \(Suz\), \(O'N\).
\item
  \(L_2(q)\), \(q\ge 4\), with \(\ \omega\big(\tfrac{q-1}{d}\big)\le 2\)
  and \(\omega\big(\tfrac{q+1}{d}\big)\le 2\), where \(d=(2,q-1)\).
\item
  \(L_3(q)\) with
  \(\ \omega\big(\tfrac{q-1}{d}\big)\le 1,\ \ \omega\big(\tfrac{q^2-1}{d}\big)\le 2,\ \ \omega\big(\tfrac{q^2+q+1}{d}\big)\le 2\),
  where \(d=(3,q-1)\).
\item
  \(U_3(q)\), \(q\ge 3\), with
  \(\ \omega\big(\tfrac{q+1}{d}\big)\le 1,\ \ \omega\big(\tfrac{q^2-1}{d}\big)\le 2,\ \ \omega\big(\tfrac{q^2-q+1}{d}\big)\le 2\),
  where \(d=(3,q+1)\).
\item
  \(Sz(q)={}^2B_2(q)\), \(q=2^{2m+1}\ge 8\), with \(\omega(q-1)\le 2\)
  and \(\omega(q\pm\sqrt{2q}+1)\le 2\).
\item
  \({}^2G_2(q)\), \(q=3^{2m+1}\ge 27\), with \(\omega(q-1)\le 2\),
  \(\omega(q+1)\le 2\) and \(\omega(q\pm\sqrt{3q}+1)\le 2\).
\item
  \textbf{Eighteen further groups of Lie type:} \(L_4(3)\), \(L_5(2)\),
  \(U_4(2)\cong PSp_4(3)\), \(U_4(3)\), \(U_5(2)\), \(U_6(2)\),
  \(PSp_4(4)\), \(PSp_4(8)\), \(PSp_6(2)\), \(\Omega_7(3)\),
  \(P\Omega_8^+(2)\), \(P\Omega_8^+(3)\), \(G_2(3)\), \(G_2(4)\),
  \(G_2(8)\), \({}^3D_4(2)\), \({}^2F_4(2)'\), \({}^2F_4(8)\).
\end{enumerate}

\end{answerbox}

(Isomorphisms such as \(A_5\cong L_2(4)\cong L_2(5)\),
\(A_6\cong L_2(9)\), \(A_8\cong L_4(2)\), \(L_2(7)\cong L_3(2)\),
\(U_3(3)\cong G_2(2)'\), \(L_2(8)\cong{}^2G_2(3)'\) cause harmless
repetitions.)

All groups of Lie type not covered by items 3--8 fail (P2); in
particular every classical group whose natural module has dimension
\(\ge 9\), every group of Lie type of BN-rank \(\ge 2\) over a field
with \(q\notin\{2,3,4,8\}\) other than \(L_3(q)\), and every group of
type \(F_4, E_6, {}^2E_6, E_7, E_8\) fails (P2).

Members of the families with small \(q\) (computed from the criteria and
cross-checked by an independent method, §10):

\begin{itemize}
\tightlist
\item
  \(L_2(q)\), \(q\le 1000\) (86 values in \(4\le q\le 1000\)): \(q=4\),
  \(5\), \(7\), \(8\), \(9\), \(11\), \(13\), \(16\), \(17\), \(19\),
  \(23\), \(25\), \(27\), \(29\), \(31\), \(32\), \(37\), \(41\),
  \(43\), \(47\), \(49\), \(53\), \(64\), \(67\), \(71\), \(73\),
  \(79\), \(81\), \(89\), \(97\), \(101\), \(\dots\);
\item
  \(L_3(q)\), \(q\le 200\): \(q=2\), \(3\), \(4\), \(5\), \(7\), \(8\),
  \(9\), \(13\), \(16\), \(17\), \(25\), \(49\), \(97\), \(193\);
\item
  \(U_3(q)\), \(q\le 200\): \(q=3\), \(4\), \(5\), \(7\), \(8\), \(11\),
  \(23\), \(32\), \(47\), \(128\);
\item
  \(Sz(q)\): \(q=8,32,128,512,2048,8192\) satisfy the condition,
  \(q=2^{15}\) does not (\(2^{15}-1=7\cdot31\cdot151\));
\item
  \({}^2G_2(q)\): \(q=27,243,2187\) satisfy the condition, \(q=3^9\)
  does not (\(3^9-1=2\cdot 13\cdot 757\)).
\end{itemize}

Whether the families 3--7 are infinite is a question of analytic number
theory (simultaneous ``few prime factors'' conditions on values of
polynomials) which we do not address. In contrast with the one-prime
(EPPO) case, where Suzuki's list is finite, the answer here is
necessarily given by arithmetic conditions.

\subsection{1. Reduction to simple
groups}\label{reduction-to-simple-groups}

\textbf{Lemma 1.1.} (P2) is inherited by subgroups and by quotient
groups, hence by all sections. Consequently, the non-abelian composition
factors of finite (P2)-groups are exactly the non-abelian simple
(P2)-groups.

\emph{Proof.} Subgroups: obvious. Quotients: let \(N\trianglelefteq H\)
and \(hN\in H/N\) of order \(m\). Since \((hN)^{|h|}=h^{|h|}N=N\), \(m\)
divides \(|h|\), so the prime divisors of \(m\) are among those of
\(|h|\) and \(\omega(m)\le\omega(|h|)\le 2\). A composition factor
\(H/N\) of \(G\) (\(N\trianglelefteq H\le G\)) is a section of \(G\).
Conversely a simple (P2)-group \(T\) is a composition factor of the
(P2)-group \(T\). \(\square\)

So the problem is equivalent to: \emph{classify the finite non-abelian
simple groups \(T\) with (P2).}

\subsection{2. Elementary lemmas}\label{elementary-lemmas}

\textbf{Lemma 2.1.} A finite abelian group \(A\) has an element whose
order has \(\ge 3\) prime divisors iff \(\omega(|A|)\ge 3\).
\emph{Proof.} If \(r_1,r_2,r_3\) are distinct primes dividing \(|A|\),
take elements \(a_i\) of order \(r_i\) (Cauchy); \(a_1a_2a_3\) has order
\(r_1r_2r_3\). The converse is trivial. \(\square\)

\textbf{Lemma 2.2 (number theory).} For a prime power \(q\):
\(\ \omega(q-1)\le 1\) and \(\omega(q+1)\le 1\) \(\iff\)
\(q\in\{2,3,4,8\}\). \emph{Proof.} If \(q\) is odd, \(q-1\) and \(q+1\)
are even prime powers, hence powers of \(2\) differing by \(2\), so
\(\{q-1,q+1\}=\{2,4\}\), \(q=3\). Let \(q=2^f\). If \(2^f-1=r^a\) with
\(f\ge 2\), \(a\ge 2\): for \(a\) odd, \(r^a+1=(r+1)(r^{a-1}-\dots+1)\)
with the second factor odd and \(>1\), impossible for a power of \(2\);
for \(a\) even, \(r^a+1\equiv 2\pmod 4\), so \(2^f=2\), contradiction.
Hence \(2^f-1\) is \(1\) or a prime, so \(f=1\) or \(f\) is prime. If
\(2^f+1=s^b\) with \(b\ge 2\): for \(b\) odd,
\(s^b-1=(s-1)(s^{b-1}+\dots+1)\) with odd second factor \(>1\),
impossible; for \(b\) even, \((s^{b/2}-1)(s^{b/2}+1)=2^f\) gives
\(s^{b/2}=3\), \(f=3\). Otherwise \(2^f+1\) is a Fermat prime and \(f\)
is a power of \(2\). So \(f\in\{1,2,3\}\), i.e.~\(q\in\{2,4,8\}\); and
\(q=2,3,4,8\) indeed satisfy the condition. \(\square\)

\textbf{Lemma 2.3.} Let \(G\) be a finite group, \(p\) a prime,
\(Z\le Z(G)\) a \(p'\)-group, and \(g=su\) with \(s\) a \(p'\)-element,
\(u\) a \(p\)-element, \(su=us\). Then \(|gZ|=|sZ|\cdot|u|\).
\emph{Proof.} \(g^k\in Z\) iff \(s^ku^k\in Z\); since the decomposition
of an element into commuting \(p'\)- and \(p\)-parts is unique and \(Z\)
consists of \(p'\)-elements, this holds iff \(u^k=1\) and \(s^k\in Z\).
\(\square\)

\textbf{Lemma 2.4.} Let \(T=G/Z\) with \(G\) a finite group of Lie type
in characteristic \(p\), \(Z\le Z(G)\) a \(p'\)-group. If a
\(p\)-element \(uZ\) centralises \(sZ\) (\(s\) semisimple), then \(u\)
centralises \(s\). \emph{Proof.} \(usu^{-1}=sz\) with \(z\in Z\); then
\(u^isu^{-i}=sz^i\), and \(u^{p^k}=1\) gives \(z^{p^k}=1\), so \(z=1\).
\(\square\)

\textbf{Lemma 2.5 (monotonicity in classical groups).} Let \(X(V)\) be
one of \(SL(V)\), \(SU(V)\), \(Sp(V)\), \(\Omega(V)\) and \(W\le V\) a
non-degenerate (any subspace for \(SL\)) subspace with a complement
\(W'\) (\(W'=W^\perp\) in the form cases), \(W'\neq 0\). Then
\(X(W)\times 1_{W'}\le X(V)\) and \((X(W)\times 1_{W'})\cap Z(X(V))=1\)
(scalars act as scalars on \(W'\)). Hence the simple group
\(X(V)/Z(X(V))\) contains a subgroup isomorphic to \(X(W)\), which maps
onto \(X(W)/Z(X(W))\). In particular \textbf{if \(X(W)/Z(X(W))\) fails
(P2), so does \(X(V)/Z(X(V))\)} (lift a bad element to \(X(W)\), Lemma
1.1). (For \(\Omega\): the spinor norm, resp. the Dickson invariant for
\(q\) even, and the determinant are multiplicative over orthogonal sums,
so \(\Omega(W)\times 1\le\Omega(V)\).) \(\square\)

\subsection{3. Alternating groups}\label{alternating-groups}

\textbf{Proposition 3.1.} \(A_n\) (\(n\ge5\)) has (P2) iff \(n\le 11\).
\emph{Proof.}
\((1\,2\,3\,4\,5)(6\,7\,8)(9\,10)(11\,12)\in A_{12}\le A_n\)
(\(n\ge12\)) has order \(30\). Let \(n\le 11\) and \(x\in A_n\) with
\(r_1r_2r_3\mid |x|\) for primes \(r_1<r_2<r_3\). If no cycle length of
\(x\) is divisible by two of the \(r_i\), the cycles cover at least
\(r_1+r_2+r_3\) points; this is \(\ge 12\) unless \(\{r_i\}=\{2,3,5\}\);
but then an even cycle is needed and, \(x\) being even, a second even
cycle: at least \(2+2+3+5=12\) points. If a cycle has length divisible
by \(r_ir_j\), it has length \(\ge 6\), and covering the third prime
gives either \(6+5=11\) (a \(6\)-cycle and a \(5\)-cycle: an odd
permutation, excluded), or at least \(10+3=13\), \(6+7=13\), \(15\)
points. Contradiction. (Also confirmed by the character tables of
\(A_5,\dots,A_{19}\), §10.) \(\square\)

\subsection{4. Sporadic groups}\label{sporadic-groups}

\textbf{Proposition 4.1.} The sporadic groups with (P2) are exactly
\(M_{11}\), \(M_{12}\), \(M_{22}\), \(M_{23}\), \(M_{24}\), \(J_1\),
\(J_2\), \(J_3\), \(HS\), \(He\), \(Ru\), \(Suz\), \(O'N\).
\emph{Proof.} Inspection of the element orders (ATLAS; GAP
\texttt{OrdersClassRepresentatives} of the 26 tables in the Character
Table Library). The other thirteen groups \(Co_1\), \(Co_2\), \(Co_3\),
\(Fi_{22}\), \(Fi_{23}\), \(Fi_{24}'\), \(HN\), \(Ly\), \(Th\), \(J_4\),
\(McL\), \(B\), \(M\) all contain elements of order \(30\). (Script
\texttt{scan\_ctbllib.g}, output \texttt{scan\_ctbllib.out}.)
\(\square\)

\subsection{5. Groups of Lie type: general facts and a decision
procedure}\label{groups-of-lie-type-general-facts-and-a-decision-procedure}

Let \(\mathbf G\) be a simple simply connected algebraic group over
\(\overline{\mathbb F}_p\), \(F\) a Steinberg endomorphism,
\(G=\mathbf G^F\), \(Z=Z(G)\). Apart from \(Sp_4(2)'\cong A_6\),
\(G_2(2)'\cong U_3(3)\), \({}^2G_2(3)'\cong L_2(8)\) and \({}^2F_4(2)'\)
(treated by their tables) and the solvable groups
\(SL_2(2),SL_2(3),SU_3(2),Sz(2)\), every simple group of Lie type is
\(T=G/Z\) for such \(G\). We use the standard facts (Carter {[}7, §§1.4,
3.3, 3.5{]}, Springer--Steinberg {[}9{]}, Digne--Michel {[}8, Ch. 3{]}):

\begin{itemize}
\tightlist
\item
  \begin{enumerate}
  \def\labelenumi{(\Alph{enumi})}
  \setcounter{enumi}{9}
  \tightlist
  \item
    every \(g\in G\) is \(g=su=us\) with \(s\) semisimple (\(p'\)) and
    \(u\) unipotent (\(p\)-element);
  \end{enumerate}
\item
  \begin{enumerate}
  \def\labelenumi{(\Alph{enumi})}
  \setcounter{enumi}{19}
  \tightlist
  \item
    every semisimple \(s\in G\) lies in an \(F\)-stable maximal torus;
    the \(G\)-classes of \(F\)-stable maximal tori correspond to
    \(F\)-classes of the Weyl group \(W\), and
    \(\mathbf T_w^F\cong Y/(F_w-1)Y\) where \(Y\) is the cocharacter
    lattice (= coroot lattice, \(\mathbf G\) simply connected) and
    \(F_w\) the action of \(Fw\) on \(Y\); equivalently
    \(\mathbf T_w^F=(F_w-1)^{-1}Y/Y\subset Y\otimes\mathbb Q/\mathbb Z\),
    and for \(s\leftrightarrow v\), a root \(\alpha\) satisfies
    \(\alpha(s)=1\iff\langle\alpha,v\rangle\in\mathbb Z\);
  \end{enumerate}
\item
  \begin{enumerate}
  \def\labelenumi{(\Alph{enumi})}
  \setcounter{enumi}{18}
  \tightlist
  \item
    (Steinberg) \(C_{\mathbf G}(s)\) is connected reductive, generated
    by \(\mathbf T\) and the root subgroups \(U_\alpha\) with
    \(\alpha\in\Phi_s:=\{\alpha:\alpha(s)=1\}\);
    \(s\in Z(C_{\mathbf G}(s))\); and \(C_G(s)=C_{\mathbf G}(s)^F\)
    contains a non-trivial unipotent element iff
    \(\Phi_s\ne\varnothing\).
  \end{enumerate}
\end{itemize}

By Lemma 2.3 (Lemma 2.4 only confirms that working with lifts in \(G\)
loses nothing): \textbf{\(T\) has (P2) iff (a) \(\omega(|sZ|)\le 2\) for
all semisimple \(s\in G\), and (b) \(\omega(|sZ|)\le 1\) for every
semisimple \(s\) with \(\Phi_s\neq\varnothing\).} By Lemma 2.1, (a) is
equivalent to \(\omega(|\mathbf T_w^F/Z|)\le 2\) for all \(w\). Here
\(|sZ|\) is the order of \(v\) modulo the coweight lattice, i.e.~the lcm
of the denominators of \(\langle\alpha_i,v\rangle\) for the simple roots
\(\alpha_i\).

This gives a finite decision procedure for each pair (type, \(q\)),
implemented in \texttt{torus\_method.g} (it enumerates \(\mathbf T_w^F\)
for one \(w\) in each \(F\)-class; for the Suzuki--Ree groups \(F\) acts
on \(Y\) as \(p^m\tau w\) with \(\tau(\alpha^\vee)=\alpha_\rho^\vee\) or
\(p\,\alpha_\rho^\vee\) according as \(\alpha^\vee\) is short or long,
\(q=p^{2m+1}\)). As a control, the program also checks Steinberg's
theorem that the number of \(F\)-stable maximal tori is \(|G|_p^2\),
i.e.~\(\frac{1}{|W|}\sum_{w}|\mathbf T_w^F|^{-1}=|G|_p^2/|G|\); in all
55 test cases it reproduced \(|G|\) exactly. The procedure was validated
against the character tables of 54 simple groups of Lie type (all
verdicts agree, §10). It is used as the \emph{main} argument only for
\(G_2(8)\) (for which a hand proof is also given in §8.3), and otherwise
as an independent check.

\subsection{6. The rank-one-type
families}\label{the-rank-one-type-families}

\textbf{Proposition 6.1 (\(L_2(q)\)).} \(L_2(q)\), \(q\ge4\), has (P2)
iff \(\omega(\frac{q-1}{d}),\omega(\frac{q+1}{d})\le 2\), \(d=(2,q-1)\).
\emph{Proof.} By Dickson's theorem {[}2, II.8.3--8.5, 8.27{]} every
element of \(L_2(q)\) lies in a Sylow \(p\)-subgroup (elementary
abelian) or in a cyclic subgroup of order \((q-1)/d\) or \((q+1)/d\),
and such cyclic subgroups exist. So the element orders are \(p\) and the
divisors of \((q\pm1)/d\). \(\square\)

\textbf{Proposition 6.2 (\(Sz(q)\)).} \(Sz(q)\) has (P2) iff
\(\omega(q-1),\omega(q\pm\sqrt{2q}+1)\le2\). \emph{Proof.} By Suzuki
{[}3{]} (see also Shi {[}5{]}) the element orders of \(Sz(q)\) are
exactly the divisors of \(4\), \(q-1\), \(q+\sqrt{2q}+1\),
\(q-\sqrt{2q}+1\), and cyclic subgroups of these orders exist.
\(\square\)

\textbf{Proposition 6.3 (\({}^2G_2(q)\)).} \({}^2G_2(q)\),
\(q=3^{2m+1}\ge27\), has (P2) iff
\(\omega(q-1),\omega(q+1),\omega(q\pm\sqrt{3q}+1)\le 2\). \emph{Proof.}
By Brandl--Shi {[}4, Lemma 4{]} (see also Ward {[}6{]}) the set of
element orders of \({}^2G_2(q)\) is the set of divisors of
\(6,9,q-1,\frac{q+1}{2},q\pm\sqrt{3q}+1\); since \(4\mid q+1\),
\(\omega(\frac{q+1}{2})=\omega(q+1)\). \(\square\)

\textbf{Proposition 6.4 (\(L_3(q)\) and \(U_3(q)\)).} Let
\(\varepsilon=+1\) (\(L_3(q)\)) or \(\varepsilon=-1\) (\(U_3(q)\),
\(q\ge3\)), \(d=(3,q-\varepsilon)\). Then \(L_3^\varepsilon(q)\) has
(P2) iff
\[\omega\Big(\frac{q-\varepsilon}{d}\Big)\le 1,\qquad \omega\Big(\frac{q^2-1}{d}\Big)\le 2,\qquad \omega\Big(\frac{q^2+\varepsilon q+1}{d}\Big)\le 2 .\]

\emph{Proof.} Let \(S=SL_3(q)\) acting on \(V=\mathbb F_q^3\), resp.
\(S=SU_3(q)\) acting on \(V=\mathbb F_{q^2}^3\) with a non-degenerate
hermitian form; \(Z=Z(S)\), \(|Z|=d\), \(T=S/Z\). Let
\(\mu_{q-\varepsilon}\le\overline{\mathbb F}_p^{\,*}\) be the group of
\((q-\varepsilon)\)-th roots of unity.

\emph{Semisimple elements.} Define abelian subgroups of \(S\): \(A_1\) =
the diagonal subgroup (\(\varepsilon=-1\): diagonal w.r.t. an
orthonormal basis), \(|A_1|=(q-\varepsilon)^2\); \(A_2\): for
\(\varepsilon=+1\), \(\{x\oplus N(x)^{-1}\}\) on
\(\mathbb F_{q^2}\oplus\mathbb F_q\), \(x\in\mathbb F_{q^2}^*\); for
\(\varepsilon=-1\), \(\{x\) on \(e\), \(x^{-q}\) on \(f\), \(x^{q-1}\)
on \(w\}\) with \((e,f)\) a hyperbolic pair and \(w\) anisotropic,
\(x\in\mathbb F_{q^2}^*\); \(|A_2|=q^2-1\); \(A_3\): the elements of
determinant \(1\) of the ``Singer'' torus \(\mathbb F_{q^3}^*\) (resp.
of \(\mu_{q^3+1}\le\mathbb F_{q^6}^*\) acting on
\(\mathbb F_{q^6}\cong\mathbb F_{q^2}^3\) with the trace hermitian
form); \(|A_3|=q^2+\varepsilon q+1\) (for \(\varepsilon=-1\):
\(\det x=x^{1+q^2+q^4}=x^{q^2-q+1}\) on \(\mu_{q^3+1}\)). Each \(A_i\)
contains \(Z\). Let \(s\in S\) be semisimple; we show that \(s\) lies in
an abelian subgroup \(A\le S\) with \(Z\le A\) and
\(|A|\in\{|A_1|,|A_2|,|A_3|\}\). For \(\varepsilon=+1\), according to
the factorisation of the characteristic polynomial into irreducibles
over \(\mathbb F_q\) (degrees \(1+1+1\), \(1+2\), \(3\)), \(s\) is
diagonalisable over \(\mathbb F_q\), or acts as multiplication by some
\(x\in\mathbb F_{q^2}^*\) on a 2-dimensional summand
\(\cong\mathbb F_{q^2}\) and by \(N(x)^{-1}\) on a complementary line,
or generates \(\mathbb F_q[s]\cong\mathbb F_{q^3}\); so \(s\) lies in a
\(GL_3(q)\)-conjugate of \(A_1\), \(A_2\) or \(A_3\). For
\(\varepsilon=-1\), decompose \(V\) into the primary components \(V_f\)
of \(s\) (\(f\) irreducible over \(\mathbb F_{q^2}\)); \(V_f\perp V_g\)
unless \(g=f^*\) (the polynomial whose roots are the \(\lambda^{-q}\)).
A self-dual \(f\) of degree \(2\) is impossible (it would force
\(\lambda^{q+1}=1\), i.e.~\(\lambda\in\mathbb F_{q^2}\)), so there are
three cases. (i) All \(f\) are self-dual of degree \(1\): the
eigenspaces are non-degenerate and orthogonal, so \(s\) is diagonal in
an orthonormal basis, i.e.~lies in a conjugate of \(A_1\). (ii)
\(V=V_f\oplus V_{f^*}\perp\langle w\rangle\) with \(\deg f=1\),
\(f\ne f^*\): choose \(e\in V_f\), \(f'\in V_{f^*}\) with \((e,f')=1\);
then \(s=\mathrm{diag}(x,x^{-q},x^{q-1})\) in the basis \(e,f',w\) (the
last entry is forced by \(\det s=1\)), so \(s\) lies in a conjugate of
\(A_2\). (iii) \(f\) self-dual of degree \(3\): then
\(K=\mathbb F_{q^2}[s]\cong\mathbb F_{q^6}\), the adjoint involution of
the form restricts to the automorphism \(y\mapsto y^{q^3}\) of \(K\)
(the only one of order \(2\) extending \(y\mapsto y^q\) on
\(\mathbb F_{q^2}\)), so \(A:=\{y\in K^*: y^{q^3+1}=1,\ \det y=1\}\) is
an abelian subgroup of \(S\) containing \(s\) and \(Z\), of order
\(q^2-q+1\). In all cases a \(GL_3\)- or \(GU_3\)-conjugate of \(A_i\)
is again an abelian subgroup of \(S\) of order \(|A_i|\) containing
\(Z\). Hence the semisimple elements of \(T\) are exactly the elements
of the subgroups \(A/Z\), where \(A\ge Z\) runs over abelian subgroups
of \(S\) of orders \(|A_1|,|A_2|,|A_3|\) (the \(A_i\) themselves among
them), and by Lemma 2.1 they satisfy (P2) iff
\(\omega((q-\varepsilon)^2/d)=\omega(q-\varepsilon)\le2\),
\(\omega((q^2-1)/d)\le2\), \(\omega((q^2+\varepsilon q+1)/d)\le 2\).

\emph{Elements with non-trivial unipotent part.} Let \(g=su\), \(u\ne1\)
unipotent, \(su=us\). If \(s\) had three distinct eigenvalues its
centraliser in \(GL_3(\overline{\mathbb F}_p)\) would be a torus, so
\(s\) has eigenvalues \(\lambda,\lambda,\lambda^{-2}\). If
\(\lambda^3=1\), then \(s\in Z\) and \(gZ=uZ\) has \(p\)-power order,
which is harmless; so let \(\lambda^3\ne1\). The double eigenvalue is
fixed by \(x\mapsto x^q\) (resp. \(x\mapsto x^{-q}\)), so
\(\lambda\in\mu_{q-\varepsilon}\). \(u\) preserves the \(2\)-dimensional
\(\lambda\)-eigenspace and acts trivially on the
\(\lambda^{-2}\)-eigenline, so \(|u|=p\). Moreover \(s^k\in Z\) iff
\(\lambda^{3k}=1\), so \(|sZ|=|\lambda^3|\), and by Lemma 2.3
\(|gZ|=p\,|\lambda^3|\). As \(\lambda\) runs over
\(\mu_{q-\varepsilon}\), \(\lambda^3\) runs over a cyclic group of order
\((q-\varepsilon)/d\); conversely, for every
\(\lambda\in\mu_{q-\varepsilon}\),
\(s=\mathrm{diag}(\lambda,\lambda,\lambda^{-2})\) (orthonormal basis if
\(\varepsilon=-1\)) lies in \(S\) and centralises the non-trivial
unipotent elements of \(SL\) (resp. \(SU\)) of its
\(\lambda\)-eigenspace. Hence these elements have (P2) iff
\(\omega((q-\varepsilon)/d)\le1\); this also implies
\(\omega(q-\varepsilon)\le2\). \(\square\)

(Consistent with Buturlakin's description of the spectra of
\(L_3^\varepsilon(q)\) {[}10{]}; verified for all \(q\le 200\) by the
procedure of §5, §10.)

\subsection{\texorpdfstring{7. All other groups of Lie type: reduction
to
\(q\in\{2,3,4,8\}\)}{7. All other groups of Lie type: reduction to q\textbackslash in\textbackslash\{2,3,4,8\textbackslash\}}}\label{all-other-groups-of-lie-type-reduction-to-qin2348}

\textbf{Proposition 7.1.} Let \(T\) be a simple group of Lie type over
\(\mathbb F_q\) not of type \(A_1\), \(A_2\), \({}^2A_2\), \({}^2B_2\),
\({}^2G_2\). If \(T\) has (P2) then \(q\in\{2,3,4,8\}\), and \(T\) is
not of type \(F_4,E_6,{}^2E_6,E_7,E_8\).

\emph{Proof.} We exhibit elements of orders \(p(q-1)\) and \(p(q+1)\);
then (P2) forces \(\omega(q\pm1)\le1\) and Lemma 2.2 gives
\(q\in\{2,3,4,8\}\). Recall \(SL_2(q)\) has cyclic subgroups of orders
\(q-1\) and \(q+1\).

\begin{itemize}
\tightlist
\item
  \(L_n(q)\), \(n\ge4\); \(U_n(q)\), \(n\ge4\); \(PSp_{2n}(q)\),
  \(n\ge2\): the natural module contains \(W_1\oplus W_2\ (\oplus W_3)\)
  with \(W_1,W_2\) two-dimensional (non-degenerate in the
  unitary/symplectic case) and \(SL(W_i)=SU(W_i)=Sp(W_i)\cong SL_2(q)\).
  Take \(u\ne1\) unipotent in \(SL(W_1)\) and \(y\in SL(W_2)\) of order
  \(q\pm1\), \(g=u\oplus y\oplus 1\). If \(g^k\) is scalar \(\lambda\),
  then \(u^k=\lambda\cdot 1_{W_1}\) is unipotent, so \(\lambda=1\);
  hence \(|gZ|=p\,(q\pm1)\).
\item
  \(\Omega_{2n+1}(q)\), \(q\) odd, \(n\ge3\), and
  \(P\Omega^\pm_{2n}(q)\), \(n\ge4\), \(q\) odd: by Lemma 2.5 they
  contain \(\Omega_5(q)\cong PSp_4(q)\), which has such elements.
  \(P\Omega^\pm_{2n}(q)\), \(q\) even, \(n\ge4\): contains
  \(Sp_{2n-2}(q)\ge Sp_4(q)\) (stabiliser of a non-singular vector,
  {[}16, 4.1.7{]}).
\item
  \(G_2(q)\): contains \(SL_3(q)\) and \(SU_3(q)\) ({[}11{]}, {[}12{]});
  in them \(s=\mathrm{diag}(\lambda,\lambda,\lambda^{-2})\) with
  \(|\lambda|=q\mp1\), times a commuting unipotent element of order
  \(p\) of \(SL\) (resp. \(SU\)) of the \(\lambda\)-eigenspace, has
  order \(p(q\mp1)\); these subgroups are contained in \(G_2(q)\)
  faithfully.
\item
  \({}^3D_4(q)\ge G_2(q)\) {[}13{]}; \({}^2F_4(q)\ge Sp_4(q)\),
  \(q=2^{2m+1}\) {[}14{]}.
\item
  \(F_4(q)\ge {}^3D_4(q).3\) {[}15{]}; by the above and §8
  (\({}^3D_4(q)\) fails (P2) for every \(q\ne2\)) and the table of
  \(F_4(2)\) (element of order \(30\)), \(F_4(q)\) fails (P2) for all
  \(q\). {[}Alternatively \(F_4(q)\ge B_4(q)\), and \(\Omega_9(q)\),
  \(Sp_8(q)\) fail.{]} Finally \(F_4(q)\) embeds in \(E_6(q)\),
  \({}^2E_6(q)\) (fixed points of a graph automorphism of \(E_6^{sc}\)),
  and in \(E_7(q)\), \(E_8(q)\) (via the Levi subgroups
  \(E_6^{sc}\le E_7^{sc}\), \(E_7^{sc}\le E_8\); \(F_4(q)\) is simple,
  so it meets the centre trivially). \(\square\)
\end{itemize}

\subsection{\texorpdfstring{8. The cases
\(q\in\{2,3,4,8\}\)}{8. The cases q\textbackslash in\textbackslash\{2,3,4,8\textbackslash\}}}\label{the-cases-qin2348}

\subsubsection{8.1 Classical groups}\label{classical-groups}

By Lemma 2.5, once \(X_m(q)\) fails, all \(X_n(q)\), \(n\ge m\), fail.
The base cases (GAP character tables unless stated):

\begin{center}\small\begin{tabularx}{\linewidth}{ll>{\raggedright\arraybackslash}X}\toprule
family & (P2) holds & first failure (witness order) \\
\midrule
\(L_n(2)\), \(n\ge4\) & \(L_4(2)\cong A_8\), \(L_5(2)\) & \(L_6(2)\) (30) \\
\(L_n(3)\), \(n\ge4\) & \(L_4(3)\) & \(L_5(3)\) (78) \\
\(L_n(4)\), \(n\ge4\) & --- & \(L_4(4)\) (30) \\
\(L_n(8)\), \(n\ge4\) & --- & \(L_4(8)\) (126 = \(2\cdot3^2\cdot7\); explicit matrix, \texttt{boundary\_check.g}) \\
\(U_n(2)\), \(n\ge4\) & \(U_4(2)\), \(U_5(2)\), \(U_6(2)\) & \(U_7(2)\) (30) \\
\(U_n(3)\), \(n\ge4\) & \(U_4(3)\) & \(U_5(3)\) (42) \\
\(U_n(4)\), \(n\ge4\) & --- & \(U_4(4)\) (30) \\
\(U_n(8)\), \(n\ge4\) & --- & \(U_4(8)\) (455 = \(5\cdot7\cdot13\); explicit element of \(SU_4(8)\), centre trivial) \\
\(PSp_{2n}(2)\), \(n\ge3\) & \(PSp_6(2)\) & \(PSp_8(2)\) (30) \\
\(PSp_{2n}(3)\), \(n\ge2\) & \(PSp_4(3)\cong U_4(2)\) & \(PSp_6(3)\) (30) \\
\(PSp_{2n}(4)\), \(n\ge2\) & \(PSp_4(4)\) & \(PSp_6(4)\) (30) \\
\(PSp_{2n}(8)\), \(n\ge2\) & \(PSp_4(8)\) & \(PSp_6(8)\) (455; explicit element of \(Sp_6(8)\)) \\
\(\Omega_{2n+1}(3)\), \(n\ge3\) & \(\Omega_7(3)\) & \(\Omega_9(3)\) (30) \\
\(P\Omega^+_{2n}(2)\), \(n\ge4\) & \(P\Omega^+_8(2)\) & \(P\Omega^+_{10}(2)\) (30) \\
\(P\Omega^-_{2n}(2)\), \(n\ge4\) & --- & \(P\Omega^-_8(2)\) (30) \\
\(P\Omega^+_{2n}(3)\), \(n\ge4\) & \(P\Omega^+_8(3)\) & \(P\Omega^+_{10}(3)\) (30) \\
\(P\Omega^-_{2n}(3)\), \(n\ge4\) & --- & \(P\Omega^-_8(3)\) (30) \\
\(P\Omega^\pm_{2n}(4),P\Omega^\pm_{2n}(8)\), \(n\ge4\) & --- & contain \(Sp_6(q)\), which fails \\
\bottomrule\end{tabularx}\end{center}

Monotonicity used: \(L_n\supseteq SL_{n-1}\), \(U_n\supseteq SU_{n-1}\),
\(PSp_{2n}\supseteq Sp_{2n-2}\); \(\Omega_{2n+1}(3)\) (\(n\ge5\)) and
\(P\Omega^\pm_{2n}(3)\) (\(n\ge5\)) contain \(\Omega_9(3)\);
\(P\Omega^\pm_{2n}(2)\supseteq Sp_{2n-2}(2)\supseteq Sp_8(2)\) for
\(n\ge5\); \(P\Omega^-_{2n}(q)\supseteq\Omega^-_8(q)\).
(\(\Omega_{2n+1}(q)\) for \(q\) even is \(PSp_{2n}(q)\); orthogonal
groups of dimension \(\le6\) are \(L_2,L_4,U_4,PSp_4\).) In particular,
\(PSp_4(q)\) has (P2) exactly for \(q\in\{3,4,8\}\).

\subsubsection{8.2 Exceptional groups}\label{exceptional-groups}

\begin{itemize}
\tightlist
\item
  \(G_2(q)\), \(q\in\{3,4,8\}\): all three have (P2) (\(G_2(3),G_2(4)\):
  tables; \(G_2(8)\): §8.3). (\(G_2(2)'\cong U_3(3)\).)
\item
  \({}^3D_4(q)\): \({}^3D_4(2)\) has (P2) (table); \({}^3D_4(3)\) (42),
  \({}^3D_4(4)\) (42) fail (tables); \({}^3D_4(8)\) fails: it contains
  \(SL_2(512)\times SL_2(8)\) {[}13{]}, hence an element of order
  \(513\cdot7=3^3\cdot7\cdot19\) (also found by the procedure of §5).
  Other \(q\): Proposition 7.1.
\item
  \({}^2F_4(q)\): \({}^2F_4(2)'\) and \({}^2F_4(8)\) have (P2) (tables;
  for \({}^2F_4(8)\) also §5: the eleven maximal tori have orders
  \(25\), \(35\), \(37\), \(49\), \(57\), \(63\), \(65\), \(81\),
  \(91\), \(109\), \(169\), and every semisimple \(s\ne1\) with
  \(\Phi_s\neq\varnothing\) has order dividing \(3\), \(7\), \(9\),
  \(5\) or \(13\)). For \(q=2^{2m+1}\ge32\): Proposition 7.1.
\item
  \(F_4\), \(E_6\), \({}^2E_6\), \(E_7\), \(E_8\): fail (Proposition
  7.1).
\end{itemize}

\subsubsection{\texorpdfstring{8.3 Hand proof that \(G_2(8)\) has
(P2)}{8.3 Hand proof that G\_2(8) has (P2)}}\label{hand-proof-that-g_28-has-p2}

\(G_2\) is simply connected and adjoint. (a) The \(F\)-classes of
\(W=D_{12}\) give six classes of maximal tori, with
\(|\mathbf T_w^F|=\chi_w(q)\), \(\chi_w\) the characteristic polynomial
of \(w\): \((q-1)^2,(q+1)^2,q^2-1\) (two classes), \(q^2+q+1\),
\(q^2-q+1\). For \(q=8\): \(49,81,63=3^2\cdot7,73,57=3\cdot19\), all
with \(\omega\le2\). (b) With \(a\) short and \(b\) long simple roots:
if \(\alpha_1,\dots,\alpha_k\in\Phi_s\) and
\(\beta\in\Phi\cap\sum\mathbb Z\alpha_i\) then \(\beta(s)=1\), so
\(\Phi_s=\Phi\cap\mathbb Z\Phi_s\); in particular \(\Phi_s\) is a root
subsystem. The root subsystems of \(G_2\) are, up to \(W\),
\(\varnothing\), \(A_1\), \(\tilde A_1\),
\(A_1\tilde A_1=\{\pm b,\pm(2a+b)\}\),
\(A_2=\{\pm b,\pm(3a+b),\pm(3a+2b)\}\), \(\tilde A_2\) (short roots) and
\(G_2\); \(\tilde A_2\) does not occur since the short roots span the
whole root lattice, and \(\Phi_s=G_2\) means \(s\in Z(\mathbf G)=1\).
Writing \(t\in\mathbf T\) as \((a(t),b(t))=(x,y)\) (recall
\(X(\mathbf T)=\mathbb Za\oplus\mathbb Zb\)): for \(A_1\tilde A_1\),
\(Z(C_{\mathbf G}(s))=\{y=1,\ x^2y=1\}=\{y=1,x^2=1\}\) is trivial in
characteristic \(2\); for \(A_2\) it is \(\{y=1,\ x^3y=1\}\), of order
\(3\); for \(A_1\) or \(\tilde A_1\) it is the kernel of a primitive
character, hence a one-dimensional torus \(\mathbf S\), \(F\)-stable
because \(C_{\mathbf G}(s)\) is, and \(F\) acts on
\(X(\mathbf S)\cong\mathbb Z\) as \(\pm q\), so \(|\mathbf S^F|=q\mp1\).
So every semisimple \(s\ne1\) with \(\Phi_s\ne\varnothing\) has order
dividing \(3\), \(7\) or \(9\), a prime power. By §5, \(G_2(8)\) has
(P2). \(\square\)

\subsection{9. Conclusion of the proof}\label{conclusion-of-the-proof}

Every non-abelian simple group is alternating, sporadic, or of Lie type.
§§3--4 settle the first two classes. For Lie type: the families
\(L_2,L_3,U_3,Sz,{}^2G_2\) are settled by §6 (for every \(q\)); by
Proposition 7.1 every other group of Lie type with (P2) is defined over
\(q\in\{2,3,4,8\}\) and is not of type \(F_4,E_6,{}^2E_6,E_7,E_8\); §8
lists exactly which of these have (P2): item 8 of the Main Theorem
together with \(L_4(2)\cong A_8\) and \(PSp_4(2)'\cong A_6\). With Lemma
1.1 this proves the Main Theorem. \(\blacksquare\)

\subsection{10. Computer verification (GAP 4.14, CTblLib, all scripts in
this
folder)}\label{computer-verification-gap-4.14-ctbllib-all-scripts-in-this-folder}

\begin{itemize}
\tightlist
\item
  \texttt{scan\_ctbllib.g} → \texttt{scan\_ctbllib.out}: all 154
  non-abelian simple groups with a character table in CTblLib (all 26
  sporadic groups, \(A_5\)--\(A_{19}\), many groups of Lie type
  including \(S_4(8)\) and \({}^2F_4(8)\)). 85 have (P2), 69 do not
  (with a witness element order).
\item
  \texttt{compare\_theorem\_with\_tables.py} →
  \texttt{compare\_theorem\_with\_tables.out}: the Main Theorem predicts
  the verdict of every one of these 154 tables correctly (no
  mismatches).
\item
  \texttt{torus\_method.g}: the decision procedure of §5.
\item
  \texttt{validate\_torus\_method.g} →
  \texttt{validate\_torus\_method.out}: 55 groups of Lie type; for the
  54 with a table the procedure agrees with the table, and Steinberg's
  count reproduces \(|G|\) in all cases.
\item
  \texttt{families\_check.g} → \texttt{families\_check.out}: the
  closed-form criteria of §6 agree with the procedure for \(L_2(q)\),
  \(q\le1000\) (191 values), \(L_3(q)\) and \(U_3(q)\), \(q\le200\),
  \(Sz(q)\), \(q\le2^{11}\), \({}^2G_2(q)\), \(q\le3^7\): 0 mismatches.
  \texttt{families\_check2.g} → \texttt{families\_check2.out} extends
  the \(L_2(q)\) check to \(1000<q\le2000\) (140 values, 0 mismatches).
  \texttt{families\_check3.g} → \texttt{families\_check3.out} extends
  the \(L_3(q)\), \(U_3(q)\) check to \(200<q\le400\) (37 values, 0
  mismatches). \texttt{direct\_L3U3\_check.g} →
  \texttt{direct\_L3U3\_check.out} checks the criterion of Proposition
  6.4 directly on the conjugacy classes of the actual groups \(L_3(q)\),
  \(q\le13\), and \(U_3(q)\), \(3\le q\le13\) (17 groups, 0 mismatches).
\item
  \texttt{mixed\_orders.g} → \texttt{mixed\_orders.out}: for \(G_2(8)\),
  \({}^2F_4(8)\), \(PSp_4(8)\) the torus orders and the orders of the
  semisimple \(s\ne1\) with \(\Phi_s\ne\varnothing\) (resp.
  \(\{3,7,9\}\), \(\{3,5,7,9,13\}\), \(\{3,7,9\}\)), confirming
  §8.2--8.3.
\item
  \texttt{extra\_torus\_check.g} → \texttt{extra\_torus\_check.out}: the
  procedure confirms that \(U_7(2)\), \(PSp_6(4)\), \(\Omega_9(3)\),
  \(P\Omega^\pm_{10}(2)\), \(P\Omega^\pm_{10}(3)\) fail (P2).
\item
  \texttt{direct\_classes\_check.g} →
  \texttt{direct\_classes\_check.out} (no character tables used):
  conjugacy classes of the actual groups \(L_4(3)\), \(L_5(2)\),
  \(U_4(3)\), \(U_5(2)\), \(PSp_4(4)\), \(PSp_6(2)\), \(PSp_4(8)\),
  \(P\Omega_8^+(2)\) (permutation action on projective points) give the
  same element orders, all with \(\omega\le2\).
  \texttt{direct\_classes\_check2.g} →
  \texttt{direct\_classes\_check2.out} does the same for \(\Omega_7(3)\)
  and \(P\Omega_8^+(3)\) (element orders \(1\)--\(10\), \(12\), \(13\),
  \(14\), \(15\), \(18\), \(20\) in both).
  \texttt{direct\_classes\_check3.g} →
  \texttt{direct\_classes\_check3.out}: \(U_6(2)\) (action of
  \(SU_6(2)\) on the 693 isotropic points; element orders
  \(1\)--\(12,15,18\)).
\item
  \texttt{g2\_8\_sample.g} → \texttt{g2\_8\_sample.out}: an independent
  construction of \(G_2(8)\) from the Chevalley lattice of the
  7-dimensional module (order \(4329310519296\) confirmed); 3000 random
  elements have orders in \(\{3,4,6,7,8,9,12,14,18,19,21,57,63,73\}\),
  none with three prime divisors (evidence only; the proof is §8.3).
\item
  \texttt{direct\_L2\_check.g} → \texttt{direct\_L2\_check.out}: the
  criterion of Proposition 6.1 agrees with the conjugacy classes of
  \(PSL(2,q)\) for all 42 prime powers \(4\le q\le128\).
\item
  \texttt{boundary\_check.g} → \texttt{boundary\_check.out}: \(L_4(8)\),
  \(U_4(8)\), \(PSp_6(8)\), \({}^3D_4(4)\), \({}^3D_4(8)\) fail,
  \(G_2(8)\) has (P2) (procedure); explicit elements of orders \(126\)
  in \(SL_4(8)\), \(30\) in \(SL_4(4)\), \(455\) in \(SU_4(8)\) and
  \(Sp_6(8)\), \(30\) in \(Sp_6(4)\) (centres trivial).
\end{itemize}

All scripts are plain GAP 4.14 files
(\texttt{gap\ -q\ \textless{}script\textgreater{}.g});
\texttt{torus\_method.g} is read by the other scripts.

\subsection{11. Status of the argument}\label{status-of-the-argument}

\begin{itemize}
\tightlist
\item
  \textbf{Proved by hand (modulo standard theory):} reduction (§1),
  alternating groups (§3), the families \(L_3(q)\), \(U_3(q)\) (§6.4),
  the reduction to \(q\in\{2,3,4,8\}\) and the bounds on rank (§7, §8.1
  monotonicity), \(G_2(8)\) (§8.3).
\item
  \textbf{Relies on cited classical results:} spectra of \(L_2(q)\)
  (Dickson), \(Sz(q)\) (Suzuki), \({}^2G_2(q)\) (Ward; Brandl--Shi); the
  subgroup containments \(SL_3(q),SU_3(q)\le G_2(q)\),
  \(G_2(q),\ SL_2(q^3)\circ SL_2(q)\le{}^3D_4(q)\),
  \(Sp_4(q)\le{}^2F_4(q)\), \({}^3D_4(q).3\le F_4(q)\),
  \(F_4(q)\le E_6(q),{}^2E_6(q),E_7(q),E_8(q)\); the general theory of
  maximal tori and semisimple centralisers (Steinberg,
  Springer--Steinberg, Carter).
\item
  \textbf{Relies on computation with library character tables
  (ATLAS/CTblLib):} the sporadic groups and the finitely many base cases
  in §8 (\(L_4(3)\), \(L_5(2)\), \(L_6(2)\), \(L_5(3)\), \(L_4(4)\),
  \(U_4(2)\), \(U_4(3)\), \(U_4(4)\), \(U_5(2)\), \(U_5(3)\),
  \(U_6(2)\), \(U_7(2)\), \(PSp_4(4)\), \(PSp_4(8)\), \(PSp_6(2)\),
  \(PSp_6(3)\), \(PSp_6(4)\), \(PSp_8(2)\), \(\Omega_7(3)\),
  \(\Omega_9(3)\), \(P\Omega^\pm_8(2)\), \(P\Omega^\pm_8(3)\),
  \(P\Omega^\pm_{10}(2)\), \(P\Omega^+_{10}(3)\), \(G_2(3)\),
  \(G_2(4)\), \({}^3D_4(2)\), \({}^3D_4(3)\), \({}^3D_4(4)\),
  \({}^2F_4(2)'\), \({}^2F_4(8)\), \(F_4(2)\)). Each of these verdicts
  except the one for \({}^2F_4(2)'\) (not of the form \(\mathbf G^F/Z\))
  is independently reproduced by the procedure of §5
  (\texttt{validate\_torus\_method.out},
  \texttt{extra\_torus\_check.out}, \texttt{boundary\_check.out}).
\item
  An independent referee pass (independent code and GAP scripts:
  character tables, direct conjugacy classes of \(PSL(3,q)\) for 22
  values \(q\le64\) and \(PSU(3,q)\) for \(q\le13\), an independent
  construction of \(G_2(8)\)) found no mathematical error; its minor
  wording remarks are incorporated.
\item
  Nothing in the classification is conjectural. The only open aspect is
  number-theoretic: whether each of the families 3--7 contains
  infinitely many groups.
\item
  Literature search (arXiv API, web search, October 2026) found no
  previous classification of simple groups with (P2); Problem 20.87 is
  still listed as open in arXiv:1401.0300v48.
\end{itemize}

\subsection{\texorpdfstring{Appendix A. Maximal element orders of the
groups in items 2 and 8 (from CTblLib; \texttt{spectra\_list.g} →
\texttt{spectra\_list.out})}{Appendix A. Maximal element orders of the groups in items 2 and 8 (from CTblLib; spectra\_list.g → spectra\_list.out)}}\label{appendix-a.-maximal-element-orders-of-the-groups-in-items-2-and-8-from-ctbllib-spectra_list.g-spectra_list.out}

\begin{center}\small\begin{tabular}{ll}\toprule
group & maximal element orders \(\mu(T)\) \\
\midrule
M11 & 5,6,8,11 \\
M12 & 6,8,10,11 \\
M22 & 5,6,7,8,11 \\
M23 & 6,8,11,14,15,23 \\
M24 & 8,10,11,12,14,15,21,23 \\
J1 & 6,7,10,11,15,19 \\
J2 & 7,8,10,12,15 \\
J3 & 8,9,10,12,15,17,19 \\
HS & 7,8,11,12,15,20 \\
He & 8,10,12,15,17,21,28 \\
Ru & 14,15,16,20,24,26,29 \\
Suz & 11,13,14,15,18,20,21,24 \\
ON & 11,12,15,16,19,20,28,31 \\
\bottomrule\end{tabular}\end{center}

\begin{center}\small\begin{tabular}{ll}\toprule
group & maximal element orders \(\mu(T)\) \\
\midrule
L4(3) & 8,9,12,13,20 \\
L5(2) & 8,12,14,15,21,31 \\
U4(2) & 5,9,12 \\
U4(3) & 5,7,8,9,12 \\
U5(2) & 8,11,12,15,18 \\
U6(2) & 7,8,10,11,12,15,18 \\
S4(4) & 4,6,10,15,17 \\
S4(8) & 4,14,18,63,65 \\
S6(2) & 7,8,9,10,12,15 \\
O7(3) & 8,12,13,14,15,18,20 \\
O8+(2) & 7,8,9,10,12,15 \\
O8+(3) & 8,12,13,14,15,18,20 \\
G2(3) & 7,8,9,12,13 \\
G2(4) & 8,10,12,13,15,21 \\
3D4(2) & 8,12,13,18,21,28 \\
2F4(2)' & 10,12,13,16 \\
2F4(8) & 12,16,18,20,28,35,37,52,57,63,65,91,109 \\
G2(8) & (no library table; (P2) proved in §8.3) \\
\bottomrule\end{tabular}\end{center}

\subsection{References}\label{references}

\begin{enumerate}
\def\labelenumi{\arabic{enumi}.}
\tightlist
\item
  E. I. Khukhro, V. D. Mazurov (eds.), \emph{Unsolved Problems in Group
  Theory. The Kourovka Notebook}, No.~21, arXiv:1401.0300v48 (2026),
  Problem 20.87.
\item
  B. Huppert, \emph{Endliche Gruppen I}, Springer 1967, Ch. II §8
  (Dickson's theorem on \(PSL_2(q)\)).
\item
  M. Suzuki, On a class of doubly transitive groups, \emph{Ann. of
  Math.} 75 (1962) 105--145.
\item
  R. Brandl, W. J. Shi, A characterization of finite simple groups with
  abelian Sylow 2-subgroups, \emph{Ricerche Mat.} 42 (1993) 193--198
  (Lemma 4: spectrum of \({}^2G_2(q)\)); the statement is quoted, e.g.,
  in arXiv:1606.00049.
\item
  W. J. Shi, A characterization of Suzuki's simple groups, \emph{Proc.
  Amer. Math. Soc.} 114 (1992) 589--591.
\item
  H. N. Ward, On Ree's series of simple groups, \emph{Trans. Amer. Math.
  Soc.} 121 (1966) 62--89.
\item
  R. W. Carter, \emph{Finite Groups of Lie Type: Conjugacy Classes and
  Complex Characters}, Wiley 1985.
\item
  F. Digne, J. Michel, \emph{Representations of Finite Groups of Lie
  Type}, 2nd ed., CUP 2020, Ch. 3.
\item
  T. A. Springer, R. Steinberg, Conjugacy classes, in: \emph{Seminar on
  Algebraic Groups and Related Finite Groups}, LNM 131, Springer 1970.
\item
  A. A. Buturlakin, Spectra of finite linear and unitary groups,
  \emph{Algebra and Logic} 47 (2008) 91--99.
\item
  P. B. Kleidman, The maximal subgroups of the Chevalley groups
  \(G_2(q)\) with \(q\) odd, the Ree groups \({}^2G_2(q)\), and of their
  automorphism groups, \emph{J. Algebra} 117 (1988) 30--71.
\item
  B. N. Cooperstein, Maximal subgroups of \(G_2(2^n)\), \emph{J.
  Algebra} 70 (1981) 23--36.
\item
  P. B. Kleidman, The maximal subgroups of the Steinberg triality groups
  \({}^3D_4(q)\) and of their automorphism groups, \emph{J. Algebra} 115
  (1988) 182--199.
\item
  G. Malle, The maximal subgroups of \({}^2F_4(q^2)\), \emph{J. Algebra}
  139 (1991) 52--69.
\item
  M. W. Liebeck, J. Saxl, G. M. Seitz, Subgroups of maximal rank in
  finite exceptional groups of Lie type, \emph{Proc. London Math. Soc.}
  65 (1992) 297--325.
\item
  P. Kleidman, M. Liebeck, \emph{The Subgroup Structure of the Finite
  Classical Groups}, CUP 1990.
\item
  J. H. Conway et al., \emph{ATLAS of Finite Groups}, OUP 1985; T.
  Breuer, \emph{The GAP Character Table Library} (CTblLib), GAP 4.14.
\item
  W. J. Shi, W. Yang, On finite groups with elements of prime power
  orders, arXiv:2003.09445 (one-prime case).
\end{enumerate}
\end{document}
